\chapter{Introduction to Image Processing}
\label{ch:introduction}

\section{Learning Objectives}

By the end of this chapter, you should be able to:

\begin{itemize}
    \item explain what is meant by digital image processing;
    \item distinguish image processing from computer graphics;
    \item describe a digital image as numerical data;
    \item explain why arrays are central to image processing;
    \item describe the implementation-led approach used in this book; and
    \item identify the relationship between the book, its Jupyter notebooks and
          the \pyimage{} package.
\end{itemize}

\section{Introduction}

Digital image processing is concerned with representing images as data and
applying computational operations to that data. Depending on the application,
those operations may enhance an image, transform its appearance, extract
measurements, compare it with another image, reduce its storage requirements or
prepare it for further analysis.

This book approaches image processing as both a mathematical subject and a
programming discipline. The objective is not merely to learn which library
function performs a task, but to understand enough of the underlying operation
to implement, test and evaluate it.

\begin{historicalnote}
The historical UQC146S1 material described the module as covering both the C
programming language and basic image processing, with only a limited period
spent introducing C before moving on to image-processing work
\cite{johnson_uqc146_intro}. This book keeps the same practical emphasis but
uses Python as the implementation language.
\end{historicalnote}

\section{Image Processing and Computer Graphics}

The historical course material made a useful distinction between computer
graphics and image processing: computer graphics is principally concerned with
generating images, while image processing begins with existing image data and
manipulates or analyses it. That distinction is useful for orientation, even
though modern applications often combine the two areas.

In this book, an image is therefore primarily treated as data. Displaying the
image is important, but the display is only one representation of the numeric
values held by the program.

\section{Images as Numerical Data}

At its simplest, a greyscale digital image can be represented as a rectangular
array of numbers. Each array element corresponds to a sample at a particular
location in the image. In an 8-bit greyscale representation, a value of 0 is
commonly associated with black and 255 with white, with intermediate values
representing intermediate levels of brightness.

A tiny image can therefore be constructed directly with NumPy:

\begin{listing}[htbp]
    \caption{A small greyscale image represented as a NumPy array.}
    \label{lst:first-array-python}

\begin{minted}{python}
import numpy as np

image = np.array([
    [0,   0,  40, 40],
    [0,  60, 120, 40],
    [20, 100, 220, 80],
    [0,  40,  80,  0],
], dtype=np.uint8)
\end{minted}

\end{listing}

The choice of \texttt{uint8} is significant. It limits each array element to an
unsigned 8-bit integer and therefore gives it a finite numeric range. Later
chapters will examine the consequences of overflow, underflow, clipping and
data-type conversion in detail.

\section{Images as Arrays}

The original course material emphasised that an image may be viewed logically
as a two-dimensional array while still being stored in memory or on disk as a
sequence of pixel values. That distinction remains important in modern Python:
NumPy gives us a convenient multidimensional view, while the underlying data
still has a concrete representation in memory.

For a greyscale image, a common array shape is:

\begin{equation}
    H \times W,
\end{equation}

where $H$ is the image height and $W$ is its width. A colour image may instead
use a shape such as:

\begin{equation}
    H \times W \times 3,
\end{equation}

where the final dimension contains three colour components, commonly red,
green and blue.

\section{An Implementation-Led Approach}

Many modern image-processing operations can be expressed with a single library
call. That is valuable in production software, but it can conceal the decisions
an algorithm makes about numeric range, border handling, interpolation,
normalisation or colour representation.

The approach used throughout this book is therefore:

\begin{enumerate}
    \item understand the operation conceptually;
    \item express the operation mathematically where appropriate;
    \item develop pseudocode;
    \item implement the operation directly;
    \item test it against known cases and boundary conditions;
    \item examine the result visually and numerically;
    \item move the validated implementation into the \pyimage{} package; and
    \item compare it with established library implementations.
\end{enumerate}

This reflects an important feature of the historical portfolio. The coursework
was not simply a request for output images: it required a library of functions,
documentation and a test plan, and placed explicit value on programming style
and demonstrated code\cite{johnson_uqc146_assignment}.

\section{The Companion Project}

The practical work is divided between Jupyter notebooks and a conventional
Python package. The notebooks are used to investigate algorithms, display
intermediate results and experiment with parameter values. The package contains
the reusable versions of the completed functions. A separate pytest suite
provides repeatable verification.

\begin{projectfiles}
\textbf{Companion notebook}

\texttt{notebooks/01-introduction.ipynb}

\medskip

\textbf{Reusable package}

\texttt{project/src/pyimage/}

\medskip

\textbf{Automated tests}

\texttt{project/tests/}
\end{projectfiles}

\section{From C to Python}

The historical module used C and deliberately exposed implementation details
such as numeric types, pointers, manual allocation and linkable libraries. A
Python implementation removes some of that machinery, but the underlying
questions do not disappear. Images still have dimensions and data types; binary
files still contain bytes in a defined order; numerical operations still have
ranges and precision; reusable software still needs a clear API; and algorithms
still need to be tested.

For that reason, the book will not treat Python as a way to avoid the technical
content of the original course. Instead, Python allows more time to be spent on
the image-processing algorithms while NumPy makes the representation of image
data explicit and inspectable.

\section{What the Book Will Build}

Over the following chapters the \pyimage{} package will grow to cover file I/O,
colour conversion, point operations, convolution, filtering, edge detection,
geometric operations, frame operations, compression and selected standard image
formats. The historical portfolio and sign-off material provides part of the
practical roadmap, while modern material extends the project beyond the scope
of the original module.

\section{Chapter Summary}

A digital image can be treated as structured numerical data. Image processing
operates on that data to transform, enhance, combine, compress or analyse an
image. This book develops those operations by implementing them before relying
on high-level library equivalents. The chapter notebooks support exploration;
the \pyimage{} package holds the finished implementations; and the automated
tests provide independent verification.

\section{Review Questions}

\begin{enumerate}
    \item What is the practical distinction between image processing and
          computer graphics used in this book?
    \item Why is an array a useful representation of a digital image?
    \item What information is expressed by the shape $H \times W \times 3$?
    \item Why might a direct implementation be useful even when a library
          already supplies the same operation?
    \item What different roles do the notebooks, the \pyimage{} package and the
          test suite play in the companion project?
\end{enumerate}